% Copyright (c) 2026 Eugene Joseph M. Ragasa.
% SPDX-License-Identifier: Apache-2.0
% Author: Eugene Joseph M. Ragasa
% Formal proof completed: 2026-10-27
\documentclass[11pt]{article}
\usepackage[T1]{fontenc}\usepackage{amsmath,amssymb}\usepackage[margin=1in]{geometry}\usepackage[hidelinks]{hyperref}
\title{Proof Dossier: Finite Fourier Coefficients Separate Fields}\author{Eugene Joseph M. Ragasa}\date{Formal proof completed: 2026-10-27}
\begin{document}\maketitle
\section{Identity and status}
ID \texttt{BG-FOURIER-THM-051}; status \texttt{proved}. Lean declaration:
\par\noindent{\tiny\path{ProjectKoios.BandGauge.Fourier.}}\\
{\tiny\path{exists_finiteFourierCoefficient_ne_of_ne}}.
\section{Statement and proof}
Two unequal fields over the declared finite abelian group differ in at least one Fourier
coefficient. Otherwise all coefficients agree, and exact finite reconstruction makes the
fields equal pointwise, contradicting the premise.
\section{Boundary}
Completeness is relative to the supplied finite group and its full character inventory; it
is not external mesh coverage or continuum uniqueness. Kernel axioms are \texttt{propext},
\texttt{Classical.choice}, and \texttt{Quot.sound}; no admissions. Checked with Lean
\texttt{4.35.0} and Mathlib revision
\texttt{56a005855ea17e0b0958f73ec801ae140b0a1679}. AI-assisted and kernel-checked;
no independent expert review recorded.
\end{document}
